% ПРИЛОЖЕНИЕ, без номер на раздел.
% Изходният текст НЕ се вмъква тук: адресът на хранилището върши същата работа.
\section*{ПРИЛОЖЕНИЕ}
\label{sec:appendix}
\addcontentsline{toc}{section}{ПРИЛОЖЕНИЕ}

Целият изходен текст, записаните изпълнения на измерванията в \code{results/} и този
документ са в публичното хранилище \url{https://github.com/Alex-Tsvetanov/pointforge}.
Табл.~\ref{tab:files} дава разпределението на кода; броят редове е за заглавния файл и
реализацията взети заедно. Изгражда се с \code{cmake -S . -B build -DCMAKE\_BUILD\_TYPE=Release}
и \code{cmake --build build}, след което \code{ctest --test-dir build} пуска проверките,
целта \code{demo} произвежда изходните файлове, а \code{pointforge\_bench --csv <файл>}
повтаря измерванията. Втората конфигурация се получава с
\code{-DPOINTFORGE\_NATIVE\_ARCH=ON} и не е преносима, затова е изключена по подразбиране.
Изграждането и \code{ctest} на Ubuntu са записани в \code{.github/workflows/ci.yml}; там не се
пускат санитайзери и покритие, и за тях не се претендира.

\begin{table}[H]
\centering
\caption{Модули на реализацията}
\label{tab:files}
\begin{tabular}{@{}L{4.4cm}rL{7.4cm}@{}}
\toprule
\textbf{Файл} & \textbf{Редове} & \textbf{Съдържание} \\
\midrule
\code{point\_cloud.hpp}        & 153 & облак от точки, разположен по компонента \\
\code{transform.\{hpp,cpp\}}   & 305 & матрица 3x3, сингулярно разлагане, кораво преобразувание \\
\code{kdtree.\{hpp,cpp\}}      & 637 & построяване, трите вида заявка, трите пътя на лист \\
\code{icp.\{hpp,cpp\}}         & 257 & съответствия, отхвърляне, оценяване по Кабш \\
\code{voxel\_grid.\{hpp,cpp\}} & 149 & прореждане с вокселна решетка \\
\code{cloud\_io.\{hpp,cpp\}}   & 307 & четене и запис на PLY и PCD \\
\code{image.\{hpp,cpp\}}       & 179 & двата растерни типа и преобразуванията между тях \\
\code{image\_io.\{hpp,cpp\}}   & 163 & четене и запис на PGM и PPM \\
\code{image\_ops.\{hpp,cpp\}}  & 238 & изглаждане, градиент, немаксимуми, прагове \\
\code{labeling.\{hpp,cpp\}}    & 179 & свързани компоненти и описанията им \\
\code{synthetic.\{hpp,cpp\}}   & 272 & генератор на изкуствени облаци и изображения \\
\code{timing.hpp}              &  80 & измерване на време без външна библиотека \\
\code{text.hpp}                &  37 & подравняване на изхода при кирилица \\
\code{apps/demo.cpp}           & 256 & демонстрация: една команда, видим изход \\
\code{benchmarks/bench.cpp}    & 442 & измервания, изход и във формат CSV \\
\code{tests/}                  & 1\,385 & рамка и шест групи проверки, 60 случая \\
\bottomrule
\end{tabular}
\end{table}
